\chapter{Cache Coherence, and What It Does Not Solve}
\label{ch:coherence}

Chapter~2 established that private caches admit multiple copies of a line.
Coherence is the discipline that keeps those copies in agreement.
This chapter
develops it far enough to make two points precise: what coherence guarantees, and
--- more importantly for this paper --- what it leaves entirely unconstrained.

\section{The Coherence Invariants}
\label{sec:invariants}

Following the formulation of Nagarajan et al.\ \cite{nagarajan2020primer}, a
memory system is coherent if it maintains two invariants.

\begin{description}
\item[Single-writer, multiple-reader (SWMR).] For any memory location and any
  point in logical time, either exactly one core may write the location and no
  core may read it, or any number of cores may read it and none may write.
\item[Data value.] The value of a location at the start of an epoch is the value
  written at the end of the last read--write epoch for that location.
\end{description}

The SWMR invariant is the operative one.
It partitions the lifetime of each line
into epochs, alternating between a single exclusive writer and a set of concurrent
readers, and the protocol's task is to enforce that partition.

An equivalent characterisation, and the one used in Sarangi's companion volume
\cite{sarangi2023nextgen}, is that coherence is \emph{per-location sequential
consistency} (PLSC): for each address considered in isolation, the machine behaves
as though all accesses to that address occurred in a single total order consistent
with each core's program order.
Stated this way, the relationship to consistency
becomes clear and is developed in \S\ref{sec:coh-vs-cons}.

\section{Protocol Families}
\label{sec:protocols}

\subsection{Write-Invalidate versus Write-Update}

Two families satisfy SWMR. A \emph{write-invalidate} protocol requires a core
intending to write to first invalidate all other copies.
A \emph{write-update}
protocol instead broadcasts the new value to all holders.
The prescribed text
treats both (\S12.4.6); invalidate protocols dominate in practice, because update
protocols consume interconnect bandwidth proportional to the number of sharers
regardless of whether those sharers will ever read the line again.
All protocols
considered in this paper are invalidate-based.

\subsection{MSI}

The minimal invalidate protocol assigns each cached line one of three states.

\begin{description}
\item[M (Modified).] This cache holds the only valid copy, and it differs from
  memory.
The core may read or write freely.
\item[S (Shared).] This cache holds a valid, clean copy; others may too. Reads
  are permitted, writes are not.
\item[I (Invalid).] No usable copy is present.
\end{description}

\begin{figure}[H]
\centering
\begin{tikzpicture}[
    st/.style={circle, draw, minimum size=11mm, font=\bfseries, fill=black!5},
    tr/.style={-{Latex[length=2mm]}, thick},
    ann/.style={font=\scriptsize, align=center}]

  \node[st] (M) at (0,3.2)   {M};
  \node[st] (S) at (-3,0)    {S};
  \node[st] (I) at (3,0)     {I};

  \draw[tr] (I) to[bend right=18] node[ann,above right,pos=0.5]{PrRd / BusRd} (S);
  \draw[tr] (S) to[bend right=18] node[ann,below,pos=0.5]{BusRdX / --} (I);
  \draw[tr] (I) to[bend right=25] node[ann,right,pos=0.5]{PrWr /\\BusRdX} (M);
  \draw[tr] (M) to[bend right=25] node[ann,left,pos=0.5]{BusRdX /\\Flush} (I);
  \draw[tr] (S) to[bend left=18]  node[ann,left,pos=0.5]{PrWr /\\BusRdX} (M);
  \draw[tr] (M) to[bend left=18]  node[ann,right,pos=0.5]{BusRd /\\Flush} (S);

  \draw[tr] (M) to[loop above] node[ann,above]{PrRd, PrWr / --} (M);
  \draw[tr] (S) to[loop left]  node[ann,left]{PrRd / --} (S);
\end{tikzpicture}
\caption{The MSI protocol. Transitions are labelled \emph{event / action}, where
\code{Pr} denotes a request from the local processor and \code{Bus} a request
observed on the interconnect.}
\label{fig:msi}
\end{figure}

\subsection{MESI and MOESI}

MSI is correct but wasteful.
A core that reads a line no other core holds, and
then writes it, must issue two bus transactions: a \code{BusRd} to enter S, and a
\code{BusRdX} to reach M. Since private, unshared data is the common case, this is
a frequent and avoidable cost.

MESI adds an \textbf{E (Exclusive)} state: the line is clean and no other cache
holds it.
A write from E proceeds silently to M with no bus transaction at all.
The additional requirement is a shared signal on the interconnect by which a
requesting core learns whether any other cache responded.

MOESI adds \textbf{O (Owned)}: a line that is dirty and shared, with one cache
designated owner and responsible for supplying it to requesters.
This permits
sharing of modified data without an intervening writeback to memory.
AMD
processors use MOESI; Intel uses a MESI derivative with an additional forwarding
state.
Full transition tables for all three protocols are deferred to
Appendix~F.

\section{Snooping and Directories}
\label{sec:snoop-dir}

The protocols above describe states and transitions; a separate question is how a
core's requests reach the others.

\emph{Snooping} broadcasts every request on a shared medium that all caches
observe.
It is simple and provides a natural total order on transactions --- the
order in which they win the bus --- but requires bandwidth proportional to the
number of cores, and does not scale past roughly sixteen.

\emph{Directory} protocols maintain, for each line, a record of which caches hold
copies, and send point-to-point messages only to those.
This scales, at the cost
of an extra indirection through the directory (a three-hop rather than two-hop
miss) and of storage for the sharer vectors.
All modern large multicores are
directory-based. The prescribed text treats directories in \S12.4.6; the
interconnect topologies over which the messages travel are the subject of \S12.6
and are relevant to this paper only insofar as they determine coherence latency,
a parameter of the simulator in Chapter~8.

\section{Coherence versus Consistency}
\label{sec:coh-vs-cons}

We can now state the distinction with the precision the rest of the paper
requires.

\begin{table}[H]
\centering
\small
\caption{The division of labour between coherence and consistency.}
\label{tab:coh-vs-cons}
\begin{tabular}{@{}p{3.4cm}p{5.2cm}p{5.2cm}@{}}
\toprule
 & \textbf{Coherence} & \textbf{Consistency} \\
\midrule
Scope        & One address at a time & Across different addresses \\
Guarantee    & A single total order of writes per location, visible to all cores
             & Which program-order relations between operations are preserved
               globally \\
Mechanism    & MSI / MESI / MOESI, snooping or directory
             & Store buffer discipline, fences, speculation \\
Where it acts & At the cache
             & Between the core and the cache \\
Failure mode & A core reads a stale value of a location
             & Two cores disagree about the relative order of accesses to
               different locations \\
\bottomrule
\end{tabular}
\end{table}

The critical structural fact is the fourth row.
\emph{The store buffer sits
between the core and the cache}, and therefore in front of the coherence
protocol entirely.
A store resting in the buffer has not yet been presented to
the coherence mechanism; there is nothing for the protocol to order, invalidate
or broadcast.
Coherence cannot police an operation it has not seen.

This is why the outcome of \S\ref{sec:motivating} survives a perfect coherence
implementation, and it is why the two topics --- which are frequently taught
together and are treated in adjacent subsections of the prescribed text --- must
be kept firmly apart.

\begin{example}[Coherence holds; the anomaly persists]
In Listing~\ref{lst:sb}, consider the accesses to \code{x}.
There are two: Core~0's
store and Core~1's load.
Coherence requires only that they be totally ordered and
that the load return the value written by the most recent preceding store in that
order.
Placing the load \emph{before} the store in coherence order satisfies this
requirement exactly, and yields $\texttt{r1} = 0$. The symmetric argument for
\code{y} yields $\texttt{r0} = 0$. Both coherence invariants hold throughout.
\end{example}

\section{What Coherence Does Contribute}
\label{sec:coh-contribution}

It would be wrong to conclude that coherence is irrelevant to consistency.
It
contributes two things that Chapter~4 will use.

First, it supplies the \emph{coherence order} relation $\co$, a total order on
the stores to each location, which is one of the primitive relations from which
consistency models are constructed.

Second, and more practically, the invalidation messages it generates constitute a
ready-made notification service.
When a core speculatively executes a load early,
the arrival of an invalidation for that line is precisely the signal that the
speculation may have been observable.
An optimistic implementation of sequential
consistency needs a detection mechanism, and coherence has already built one for
other reasons.
This observation, developed in Chapter~9, is what makes speculative
SC comparatively cheap.

\section{Summary}

Coherence is a solved problem, correctly implemented in shipping hardware, and it
guarantees per-location sequential consistency.
It says nothing about the ordering
of accesses to distinct locations, and cannot, because the mechanism that reorders
them operates upstream of it.
The specification of that cross-location ordering is
a separate architectural decision, and it is to the space of such decisions that
we now turn.
